\documentclass[conference]{IEEEtran}
\IEEEoverridecommandlockouts
% The preceding line is only needed to identify funding in the first footnote. If that is unneeded, please comment it out.

\usepackage{cite}
\usepackage{amsmath,amssymb,amsfonts}
\usepackage{algorithmic}
\usepackage{graphicx}
\usepackage{textcomp}
\usepackage{xcolor}
\usepackage{url}
\def\BibTeX{{\rm B\kern-.05em{\sc i\kern-.025em b}\kern-.08em
    T\kern-.1667em\lower.7ex\hbox{E}\kern-.125emX}}
\begin{document}

\title{Merging Traffic: Benchmark of Deep Reinforcement Learning Models}

\author{\IEEEauthorblockN{Connor Slutsky}
\IEEEauthorblockA{\textit{Fourth-Year Computer Science Major}\\
\textit{College of Engineering and Applied Sciences}\\
\textit{University of Cincinnati}\\
Cincinnati, OH \\
slutskcp@mail.uc.edu}
\and
\IEEEauthorblockN{2\textsuperscript{nd} Given Name Surname}
\IEEEauthorblockA{\textit{dept. name of organization (of Aff.)} \\
\textit{name of organization (of Aff.)}\\
\textit{name of organization (of Aff.)}\\
City, Country \\
email address or ORCID}
\and
\IEEEauthorblockN{3\textsuperscript{rd} Given Name Surname}
\IEEEauthorblockA{\textit{dept. name of organization (of Aff.)} \\
\textit{name of organization (of Aff.)}\\
\textit{name of organization (of Aff.)}\\
City, Country \\
email address or ORCID}
\and
\IEEEauthorblockN{4\textsuperscript{th} Given Name Surname}
\IEEEauthorblockA{\textit{dept. name of organization (of Aff.)} \\
\textit{name of organization (of Aff.)}\\
\textit{name of organization (of Aff.)}\\
City, Country \\
email address or ORCID}
\and
\IEEEauthorblockN{5\textsuperscript{th} Given Name Surname}
\IEEEauthorblockA{\textit{dept. name of organization (of Aff.)} \\
\textit{name of organization (of Aff.)}\\
\textit{name of organization (of Aff.)}\\
City, Country \\
email address or ORCID}
}

\maketitle


\section{Background}
\textit{In this section, provide a comprehensive overview of the problem or challenge your project
addresses. Include the motivation behind selecting this topic, its relevance to the field, and any
specific issues that your project seeks to resolve. Be sure to clearly define any terminology that is
specific to your area of research. This section sets the stage for the entire report, so ensure it's clear
and well-structured.}

Our project seeks to identify a Deep Reinforcement Learning (DRL) model that is best suited to 
equip a car to merge onto and off of a highway. We were seeking to explore a topic in the realm of 
automated driving, and after some deliberation and exploration settled on specifically the act of merging
onto and off of a highway. As drivers, many of us faced initial difficulty with this driving skill, so we
thought it would be intriguing to explore how well DRL models would fair. We also considered it potentially
challenging, such that it would be suitable for our Deep Learning class project. From our current research,
we have yet to find anything that explores similar questions or tasks (TEAM NOTE: MAYBE WE SHOULD DO A LITTLE 
MORE RESEARCH TO FIND IF THERE ARE ANY EXISTING PAPERS ON MERGING??????).

Our tool to simulate the highway environment will be the ``highway-env'' python package. This package is a 
modification of the popular ``gymnasium'' python package that contains tools for building environments suitable
for DRL or RL model training. ``highway-env'' contains many pre-designed driving environments (straight highway, 
racetrack, parking, etc.), and ours will be a modification of the ``merge-v1'' environment. 

Example of citation \cite{b1}


example of ellipses $\ldots$


\section{Related Studies}
\textit{This section should demonstrate a thorough understanding of the existing work in your field. You
are required to review at least five state-of-the-art models that are closely related to your project.
For each paper, include the following: Model Name and Authors: Cite the paper and mention the
model's creators. Methodology: Summarize the key techniques or algorithms used in the model.
Focus on how these models approach the problem at hand. Performance: Provide a comparison of
each model's performance based on standard benchmarks (e.g., accuracy, precision, recall, etc.)
and briefly mention any notable strengths or weaknesses. The discussion should not only
summarize but also critically analyze how these existing models have informed your own project.
Highlight how your approach builds upon or deviates from these studies}

TODO: DOUBLE CHECK IEEE CITATION FORMAT

\begin{thebibliography}{00}
\bibitem{b1} M. Kallel, J. Holgado-Alvarez, S. Tosatto, and C. D'Eramo, ``PPO+: Enhancing proximal policy optimization,'' Reinforcement Learning Journal, 2026. [Online]. Available: \url{https://rlj.cs.umass.edu/2026/papers/Paper8.pdf}
\bibitem{b2} H. Wang and M. Yang, ``ExO-PPO: an Extended Off-policy Proximal Policy Optimization Algorithm,'' School of Aeronautics and Astronautics Organization, Feb. 2026, arXiv:2602.09726. [Online]. Available: \url{https://arxiv.org/pdf/2602.09726}
\bibitem{b3} Y. Tang, M. Rowland, R. Munos, B. Avila Pires, and W. Dabney, ``Off-policy Distributional Q($\lambda$): Distributional RL Without Importance Sampling,'' 29th International Conference of AISTATS, vol. 300, 2026. [Online]. Available: \url{https://raw.githubusercontent.com/mlresearch/v300/main/assets/tang26a/tang26a.pdf}
\bibitem{b4} A. Oberai, S. Park, S. Levine, ``Reversal Q-Learning,'' University of California, Berkeley, June 2026, arXiv:2606.17551. [Online]. Availalbe: \url{https://arxiv.org/pdf/2606.17551}
\bibitem{b5} H. Lee, et al. ``SIMBA: SIMPLICITY BIAS FOR SCALING UP PARAMETERS IN DEEP REINFORCEMENT LEARNING,'' ICLR, 2025, arXiv:2410.09754. [Online]. Available: \url{https://arxiv.org/pdf/2410.09754}
\end{thebibliography}

\end{document}
